\chapter*{Presentación}
\markboth{Presentación}{}
\addcontentsline{toc}{chapter}{Presentación}
Este volumen estudia cómo descomponer una función en frecuencias y cómo reconstruirla. Comenzamos con series periódicas, donde las frecuencias son enteros; pasamos por grupos finitos y el espacio euclidiano; finalmente construimos la medida de Haar y la transformada sobre grupos localmente compactos abelianos. La convolución, la aproximación de la identidad y las representaciones unitarias enlazan estas etapas.

El punto de partida son las lecturas de \emph{Análisis Matemático 4, Licenciatura en Física y Matemáticas}, de Miguel Ángel Valencia Bucio \cite{valencia4}. Se conservan sus 81 ejercicios numerados, con todos sus incisos, dentro de seis partes y dieciséis capítulos temáticos. La lectura 24 disponible tiene otra procedencia: es una reconstrucción editorial independiente sobre representaciones de un grupo y de su álgebra de convolución, preparada ante la ausencia del original. Sus ocho ejercicios se identifican expresamente como reconstrucción \cite{reconstruccion24}. Diez ejercicios adicionales desarrollan los complementos de esta edición. Los 99 ejercicios tienen solución.

La autoría de las lecturas de base se reconoce expresamente. La organización editorial, las explicaciones ampliadas, las soluciones, los complementos y las rectificaciones constituyen esta edición de Eder Said Martinez Ramirez. No se presenta la reconstrucción de la lectura 24 como texto recuperado de Valencia Bucio ni como prueba de lo que contenía la lectura faltante.

\section*{Apuntes y solucionario}
Los ejercicios se presentan en estos apuntes sin soluciones. Los mismos
enunciados, con la misma numeración por capítulo y sus soluciones, se reúnen
en el documento separado \emph{Solucionario} de este volumen. Las menciones
a las soluciones y a los ejercicios resueltos de esta presentación se refieren
a ese documento complementario.

\section*{Una lectura acompañada de las demostraciones}
El destinatario es un estudiante de la mitad de la licenciatura que desea avanzar hacia análisis armónico de posgrado. Esa accesibilidad exige hacer visibles los prerrequisitos. Se utilizan integración de Lebesgue, Tonelli, Fubini, convergencia dominada y espacios $L^p$ del volumen II; la geometría de Hilbert del volumen I; y dualidad, operadores, espectros, cálculo funcional y álgebras $C^*$ del volumen III. El primer recorrido, hasta el muestreo, requiere sólo una parte de estas herramientas. La teoría abstracta utiliza el conjunto de los tres volúmenes anteriores.

Las pruebas indican qué permite intercambiar una suma o una integral, cuándo se trabaja con clases de funciones iguales casi en todas partes y dónde una aproximación es uniforme, fuerte o en norma. Un resultado falso no se convierte en verdadero añadiendo cálculos: se da un contraejemplo, se formula una versión válida y se demuestra. Entre las rectificaciones aparecen factores de normalización, hipótesis de continuidad, pesos modulares y limitaciones de regularidad en grupos que no son $\sigma$-compactos.

La convergencia casi en todas partes de las sumas parciales de Fourier para toda función de $L^2$ se menciona como el teorema de Carleson, siguiendo el panorama de la lectura original. Es el único gran resultado señalado expresamente fuera del desarrollo: no se demuestra ni se utiliza para probar los resultados del libro. En cambio, la existencia y unicidad de Haar, la inversión, Plancherel, la dualidad de Pontryagin y la extensión de caracteres se desarrollan con sus argumentos.

Conviene trabajar cada ejercicio antes de consultar su solución, y comprobar en un ejemplo las constantes de cada fórmula. Las referencias internas y la tabla final permiten pasar del orden de las lecturas al orden temático. Las verificaciones propuestas dentro de la teoría se incorporan a las demostraciones, aunque no aumentan el conteo de ejercicios numerados.

\section*{Convenciones y créditos}
El producto interno es lineal en su primer argumento. En el círculo usamos $dm(t)=dt/(2\pi)$ y coeficientes $c_n(f)=\int f(t)e^{-int}dm(t)$. En $\R^d$ usamos $\widehat f(\xi)=\int f(x)e^{ix\cdot\xi}dx$ y el operador unitario $\mathcal F=(2\pi)^{-d/2}\widehat{\phantom f}$. En el grupo dual adoptamos la convención de las lecturas, $\widehat f(\chi)=\int f(x)\chi(x)dx$. El signo se especifica al identificar un grupo concreto con su dual. La medida de Haar es izquierda y $\mu(Aa)=\Delta(a)\mu(A)$; por tanto $\int f(xa)dx=\Delta(a)^{-1}\int f$. Estas decisiones fijan las constantes y los factores modulares a lo largo del volumen.

Los documentos de Egor Maximenko sobre convolución finita, núcleos aproximativos e inversión orientan los complementos concretos \cite{maximenkoConvolucionFinita,maximenkoPeriodicos,maximenkoInversion}. El desarrollo de banda limitada y series cardinales toma como referencia el apunte de Iris Paola Lozano Vite y Egor Maximenko \cite{lozanoMaximenkoPaleyWiener}. Las pruebas y soluciones se desarrollan aquí y se cotejan sus normalizaciones. La obra de Gerald B. Folland \cite{folland2016} aporta una referencia verificable para la teoría abstracta. Las referencias finales identifican documentos reales y distinguen los apuntes del curso, la reconstrucción editorial, las ampliaciones y la bibliografía de contraste.

El diseño toma como referencia visual los \emph{Apuntes de Homotopía Estable} del repositorio público de MrPaborru (\url{https://github.com/MrPaborru/Apuntes-Homotop-a-Estable}), junto con una composición tradicional de libros de matemáticas. El emblema de la ESFM procede de su sitio oficial. Su uso identifica la adscripción académica del autor; esta colección es una edición personal, no una publicación oficial ni una obra avalada por el IPN. Se conserva una paleta inspirada en los colores de la ESFM y una tipografía Latin Modern.
